% Einheit 1 von 5 – Gleichungen mit zwei Variablen und grafisches Lösen (bank/lineare-gleichungssysteme/e1.jsonl, zusammenbau v0.1)
%% TODO Zweigzeile Teil 1: was hier gelernt wird (Satz in Schülersprache aus dem Einheitstitel) – Einheit 1
\einheitenkopf[e1]{Einheit 1 von 5 · Gleichungen mit zwei Variablen und grafisches Lösen}
\zweigzeile{ab Kl. 9 (am Gymnasium ab Kl. 8) · P10}
\setcounter{aufgabe}{12}

%% TODO \verfahren{…}: Verfahrensüberschrift in Sachsprache fehlt in der Bank (2.3 b)
%% TODO Titel: Ich-kann-Satz fehlt in der Bank (Platzhalter: Kettenname) – Nr. 13
%% TODO Anweisung: ein Satz über der ersten Teilaufgabe fehlt in der Bank – Nr. 13
\begin{aufgabe}{Grafisch}
\begin{teile}
\teil $(2|6)$ – Lösung von welcher Gleichung? I: $x + y = 8$, II: $3x - y = 0$. Kreuze an.\\ \kreuz{nur von I}\\ \kreuz{nur von II}\\ \kreuz{von beiden}\\ \kreuz{von keiner}
\end{teile}
\swa{$x + y = 5$ – Lösungspaare? Ergänze die Tabelle.}{\wertetabelle{x}{y}{0,1,2,3}}
\swa{$2x + y = 9$ – Lösungspaare? Ergänze die Tabelle.}{\wertetabelle{x}{y}{0,1,2,3}}
\swa{$3x + y = 15$ – Lösungspaare? Ergänze die Tabelle.}{\wertetabelle{x}{y}{0,1,2,3}}
\swa{$x + 2y = 8$ – Lösungspaare? Ergänze die Tabelle.}{\wertetabelle{x}{y}{0,2,4,6}}
\swa{$y - x = 2$ – Lösungspaare? Ergänze die Tabelle.}{\wertetabelle{x}{y}{0,1,2,3}}
\swfrage{Was bleibt gleich, was ändert sich?}
%% TODO Darstellung für schwach fehlt in der Bank (lineare-gleichungssysteme-e1-k1-s2-v1; Abschnitt „Für schwache Schüler“)
\swz{$\text{$2x + y = 6$ – nach $y$ umstellen.}$}{}
\swa{$y - x = 1$ – zeichne die Gerade.}{\begin{ksys}[xmin=-5,xmax=5,ymin=-5,ymax=5] \end{ksys}}
\swa{Schnittpunkt? Zeichne die Geraden zu I: $y = x + 1$ und II: $y = -x + 5$ und lies den Schnittpunkt ab. ( \leerfeld | \leerfeld )}{\begin{ksys}[xmin=-5,xmax=5,ymin=-5,ymax=5] \end{ksys}}
%% TODO Darstellung für schwach fehlt in der Bank (lineare-gleichungssysteme-e1-k1-s5-v1; Abschnitt „Für schwache Schüler“)
\swz[3]{Probe: Ist $(1|5)$ die Lösung von I: $3x + y = 8$ und II: $x - y = -4$? Rechne die Probe in beiden Gleichungen.}{}
\swz{Lösung des Systems? Die Geraden I und II gehören zu den Gleichungen I und II. Lies die Lösung ab.}{\begin{ksys}[xmin=-5,xmax=5,ymin=-5,ymax=5,ablesen] \gerade{2}{-3}{I} \gerade{-1}{3}{II} \end{ksys}}
\begin{teile}
\teil I: $y = 2x + 3$, II: $y = 2x - 1$ – wie viele Lösungen? Kreuze an.\\ \kreuz{genau eine}\\ \kreuz{keine}\\ \kreuz{unendlich viele}
\end{teile}
\swa{Lösung durch Probieren? I: $x + y = 8$, II: $3x + y = 14$. Setze in I nacheinander $x = 0, 1, 2, \ldots$ ein, trage $y$ ein und prüfe II. x = \leerfeld , y = \leerfeld}{\wertetabelle{x}{y}{0,1,2,3,4,5}}
\swz[3]{Zelte: Ein Zeltlager hat $14$ Zelte, Zweierzelte und Viererzelte, mit zusammen $44$ Schlafplätzen. $x$ ist die Zahl der Zweierzelte, $y$ die Zahl der Viererzelte. Wie viele Zelte jeder Art sind es? Löse durch Probieren oder grafisch und begründe, dass es keine andere Möglichkeit gibt. \hfill (P10 2016 OS)}{}
\end{aufgabe}

%% TODO \verfahren{…}: Verfahrensüberschrift in Sachsprache fehlt in der Bank (2.3 b)
%% TODO Titel: Ich-kann-Satz fehlt in der Bank (Platzhalter: Kettenname) – Nr. 14
%% TODO Anweisung: ein Satz über der ersten Teilaufgabe fehlt in der Bank – Nr. 14
\begin{aufgabe}{Grafisch, Sek II}
\begin{teile}
\teil Schnittpunkt oder keiner? Kreuze an, wie die gezeichneten Geraden I und II liegen.\\ \kreuz{ein Schnittpunkt}\\ \kreuz{parallel}\\ \kreuz{dieselbe Gerade}
\end{teile}
\swa{Lösungsgerade? Das System I: $x - y = 2$, II: $3x - 3y = 6$ hat unendlich viele Lösungen. Zeichne alle Lösungen als Gerade und gib die Lösung mit $y = 3$ an. x = \leerfeld}{\begin{ksys}[xmin=-5,xmax=5,ymin=-5,ymax=5] \end{ksys}}
\swa{Lösungsgerade? Das System I: $2x + y = 4$, II: $-4x - 2y = -8$ hat unendlich viele Lösungen. Zeichne alle Lösungen als Gerade und gib die Lösung mit $y = 2$ an. x = \leerfeld}{\begin{ksys}[xmin=-5,xmax=5,ymin=-5,ymax=5] \end{ksys}}
\swa{Lösungsgerade? Das System I: $x + 2y = 4$, II: $2x + 4y = 8$ hat unendlich viele Lösungen. Zeichne alle Lösungen als Gerade und gib die Lösung mit $y = 1$ an. x = \leerfeld}{\begin{ksys}[xmin=-5,xmax=5,ymin=-5,ymax=5] \end{ksys}}
\swa{Lösungsgerade? Das System I: $-x + y = 1$, II: $3x - 3y = -3$ hat unendlich viele Lösungen. Zeichne alle Lösungen als Gerade und gib die Lösung mit $y = -2$ an. x = \leerfeld}{\begin{ksys}[xmin=-5,xmax=5,ymin=-5,ymax=5] \end{ksys}}
\swa{Lösungsgerade? Das System I: $3x - y = 1$, II: $-6x + 2y = -2$ hat unendlich viele Lösungen. Zeichne alle Lösungen als Gerade und gib die Lösung mit $y = 5$ an. x = \leerfeld}{\begin{ksys}[xmin=-5,xmax=5,ymin=-5,ymax=5] \end{ksys}}
\swfrage{Was bleibt gleich, was ändert sich?}
\swz[3]{Nur diese Lösung? Zeige, dass $x = 1$, $y = 3$ das System I: $2x + y = 5$, II: $x - 3y = -8$ löst, und begründe, dass es keine weitere Lösung gibt. \hfill (Abitur 2023 GK)}{}
\end{aufgabe}

%% TODO Titel: Ich-kann-Satz fehlt in der Bank (Platzhalter: Kettenname) – Nr. 15
%% TODO Anweisung: ein Satz über der ersten Teilaufgabe fehlt in der Bank – Nr. 15
\begin{aufgabe}{Grafisch – Fehler finden, Begründen}
\swz[3]{Ben liest bei I: $y = 2x - 1$ und II: $y = -x + 3$ den Schnittpunkt $(1|1)$ ab und macht keine Probe. Finde den Fehler.}{}
\swz[3]{Mia stellt um: \rechnung{3x + y &= 5 \\ y &= 3x + 5} Finde den Fehler.}{}
\swz[3]{Das System I: $-3x + y = 2$, II: $6x - 2y = -4$ hat unendlich viele Lösungen. Tim zeichnet dafür die Gerade $y = -3x + 2$. Finde den Fehler.}{}
\swfrage{Warum erfüllt der Schnittpunkt zweier Geraden beide Gleichungen des Systems? Begründe.}
\swfrage{I: $y = 3x + 1$, II: $y = 3x - 2$ – warum hat das System keine Lösung? Begründe.}
\swfrage{Warum haben zwei nicht äquivalente lineare Gleichungen mit zwei Variablen höchstens eine gemeinsame Lösung? Begründe.}
\end{aufgabe}

\uebersichtskasten{Eine Gleichung mit zwei Variablen hat viele Lösungspaare – alle zusammen liegen auf einer Geraden. Zum Zeichnen die Gleichung nach y umstellen. \\ I x + y = 7 → y = −x + 7 \quad II 2x − y = 2 → y = 2x − 2 \\ Lösung des Systems: der Punkt, der auf beiden Geraden liegt – der Schnittpunkt. Ablesen und mit der Probe in beiden Gleichungen prüfen. \\ Schnittpunkt (3 | 4): 3 + 4 = 7 (wA) \quad 6 − 4 = 2 (wA) \\ Sonderfälle: parallele Geraden (gleiche Steigung) – keine Lösung; dieselbe Gerade – unendlich viele Lösungen.}
